\section{Experimental Protocol}
\label{sec:protocol}

\subsection{Source data and source training}

The source backbone is trained \emph{only} on the de Chazal DS1 partition of the
MIT--BIH Arrhythmia Database. As published, DS1 and DS2 are \emph{record}-disjoint
rather than patient-disjoint (Section~\ref{sec:cohort}); the cohort correction
below is what makes the evaluation set genuinely patient-disjoint from DS1. The
model has 6{,}643
parameters and the architecture dimensions of Section~\ref{sec:method}. DS2
records never enter source training. On the held-out DS1 test split the frozen
source attains \macrof{} $=0.749$ over all five AAMI classes and $0.946$ over
the supported N/S/V classes, with strong V sensitivity but comparatively weak S
behavior---the regime in which patient-specific adaptation is expected to help.

\subsection{Primary and secondary cohorts}
\label{sec:cohort}

\paragraph{The de Chazal split is not patient-disjoint, and we correct it.} The
DS1/DS2 partition is the standard inter-patient protocol and we adopt it, but it
contains one subject that appears on both sides. MIT--BIH record 202's own WFDB
header states that it ``was taken from the same analog tape as record 201''
\cite{mitdb202note}, and
the two records carry identical demographics (68\,M, identical medication list).
Record \textbf{201 is in DS1} and record \textbf{202 is in DS2}. A source model
trained on DS1 has therefore already seen this subject, and evaluating adaptation
on record 202 as though it were an unseen patient measures partly memorization
rather than personalization.

We audited all 48 MIT--BIH records for this failure mode, comparing header
demographics and explicit same-tape annotations, and found exactly one such
cross-split pair: 201/202. \textbf{We exclude record 202 from the evaluation
cohort}, leaving \textbf{21 genuinely patient-disjoint records}. The source model
is unchanged, since 201 remains a legitimate DS1 training record; only the
evaluation set changes. Every primary and corrected analysis uses the 21-record
cohort. One explicitly labelled historical sensitivity result retains the earlier
22-record split and is not compared with the primary results
(Section~\ref{sec:lowbyte}). The primary grid is
$6\ \text{arms}\times21\ \text{records}\times5\ \text{seeds}=630$ adaptation
cells. We flag this because the contaminated 22-record version is what a
straightforward reading of the de Chazal protocol produces, and results computed
that way are not strictly patient-disjoint and may be optimistically biased ---
on our own data, removing record 202
raised every arm mean slightly but cost the A4--A3 comparison its statistical
significance.

\textbf{Primary cohort.} The 21 patient-disjoint DS2 records after the
\emph{post hoc} subject-identity correction form the primary cohort. Records are \emph{not}
selected using test-segment arrhythmia burden. This removes the selection bias of
the earlier exploratory panel and is the cohort behind every headline result.
Table~\ref{tab:dataset_manifest} lists the record IDs, seeds, class supports, and
split rule.

\textbf{Secondary cohort.} The 12-record arrhythmia-informative panel (records
whose test window contains $\ge 30$ arrhythmic beats) is retained only as
secondary evidence for class-specific behavior, and is clearly labeled as such
wherever it appears.

\begin{table}[H]
\centering
\small
\caption{Dataset, split, and preprocessing manifest for every database used.
MIT--BIH records are one per subject \emph{after} removing record 202, which
shares a subject with DS1 record 201 and would otherwise break patient
disjointness (Section~\ref{sec:cohort}). INCART contains multiple recordings per
subject, so its statistical unit is the subject; the full per-record inclusion
and exclusion manifests are released with the code.
\textbf{The eligibility rules are not symmetric, and the difference matters:}
primary MIT--BIH inclusion is \emph{outcome-blind} --- every corrected DS2
record is used and test-segment arrhythmia burden plays no part --- whereas the
external subsets additionally require $\ge 30$ S$+$V beats \emph{in the test
window}, so external eligibility uses future test labels and enriches those
cohorts for recordings where minority-class adaptation can be evaluated at
all.}
\label{tab:dataset_manifest}
\setlength{\tabcolsep}{5pt}
\begin{tabular}{llll}
\toprule
Field & MIT--BIH (primary) & INCART & SVDB \\
\midrule
Role & source (DS1) $+$ target (DS2) & external target & external target \\
Records evaluated & \textbf{21} (DS2) & 11 & 46 \\
Unique subjects & \textbf{21} & \textbf{6} & 46 \\
Records available & 48 & 75 & 78 \\
Lead & MLII & II & ECG1 \\
Native sampling rate & 360\,Hz & 257\,Hz & 128\,Hz \\
Resampling & none & \multicolumn{2}{l}{\texttt{resample\_poly} to 360\,Hz} \\
Beat window & \multicolumn{3}{l}{198 samples, R-peak centred (99 left / 99 right)} \\
R-peak source & \multicolumn{3}{l}{WFDB \texttt{atr} annotations} \\
Filtering & \multicolumn{3}{l}{baseline-wander removal $+$ low-pass} \\
RR normalization & \multicolumn{3}{l}{MIT--BIH DS1 statistics only, clipped at $\pm2$} \\
AAMI mapping & \multicolumn{3}{l}{N / S / V / F / Q; non-beat symbols ignored} \\
Adaptation protocol & \multicolumn{3}{l}{chronological first 500 labels, $\mathrm{gap\_beats}=1$, rest test} \\
Eligibility rule & \textbf{all corrected DS2;} & $\ge 30$ S$+$V beats & $\ge 30$ S$+$V beats \\
 & \textbf{no test-burden filter} & in the test window & in the test window \\
Eligibility uses test labels? & \textbf{no} (outcome-blind) & \textbf{yes} & \textbf{yes} \\
Seeds & \multicolumn{3}{l}{$\{42,43,44,45,46\}$} \\
\midrule
DS1 records & \multicolumn{3}{p{9.2cm}}{101 106 108 109 112 114 115 116 118 119 122 124 201 203 205 207 208 209 215 220 223 230} \\
DS2 records & \multicolumn{3}{p{9.2cm}}{100 103 105 111 113 117 121 123 200 210 212 213 214 219 221 222 228 231 232 233 234 \quad(\emph{202 excluded: same subject as DS1 record 201})} \\
\bottomrule
\end{tabular}
\end{table}


\subsection{Chronological adaptation}

Each record is split in annotation (temporal) order. The chronologically first
labeled beats form the adaptation set; a single unlabeled guard beat
($\mathrm{gap\_beats}=1$) is dropped to prevent centered-window overlap between
adaptation and test; all later beats are test-only. Early stopping uses only a
slice of the adaptation union and never the test segment. Streamed patient beats
are not counted as persistent replay storage---only the fixed source replay
memory $\mathcal{R}_0$ is. Every patient resets to $\theta_0$. The primary
run therefore contains $6\ \text{arms}\times21\ \text{records}\times5\
\text{seeds}=630$ adaptation cells (seeds $\{42,43,44,45,46\}$).

\subsection{Signal preprocessing, stated exactly}
\label{sec:preproc}

Every record passes through the same front end before beat extraction, applied to
the whole signal rather than per window. The pipeline is inherited unchanged from
the source architecture we adopt \cite{busia2025tiny}, so that our adaptation
results are attributable to the adaptation method rather than to a different
front end.
\textbf{Baseline-wander removal} uses two cascaded median filters: a 200\,ms
kernel strips the QRS and P complexes leaving baseline plus T-wave, a 600\,ms
kernel on that result strips the T-wave, and the remaining baseline estimate is
subtracted from the original signal. Kernel lengths are rounded to the nearest
odd sample count at the record's own sampling rate.
\textbf{Noise suppression} is a 4th-order Butterworth low-pass at 35\,Hz applied
\emph{zero-phase} (forward--backward, \texttt{filtfilt}), which is
non-causal: it uses future samples and could not run unmodified in a streaming
device. We keep it because it matches the source model's published front end, and
we flag it as a deployment gap rather than a result.
\textbf{Resampling} to 360\,Hz uses polyphase resampling
(\texttt{resample\_poly}) and is applied to external databases only (INCART
257\,Hz, SVDB 128\,Hz); MIT--BIH is already at 360\,Hz.

\paragraph{Non-causal preprocessing crosses the chronological boundary.} All
three of those operations --- both median kernels and the zero-phase low-pass ---
are applied to the \emph{whole record} before the chronological split, and all
three are non-causal. An adaptation-side sample within the filters' backward
reach of the split therefore depends on raw samples that belong to the test
segment. Test \emph{labels} are never involved, and no test beat enters
adaptation, but the test \emph{signal} is not strictly unseen, so the released
protocol is properly described as \textbf{offline chronological label
partitioning with non-causal whole-record preprocessing} rather than as strictly
causal chronological adaptation.

We bounded the reach empirically rather than assuming it: injecting a truncation
at a boundary and measuring how far the difference propagates backwards gives a
one-sided influence of $84$ samples at a $10^{-3}$ relative tolerance and
$120$ samples ($333$\,ms) at $10^{-5}$ --- shorter than a single RR interval.
Because a beat window spans $198$ samples, only the last one or two adaptation
beats of a record can be affected at all. Section~\ref{sec:splitfirst} reports a
split-first rerun that removes the dependency entirely and quantifies what it
changes.
\textbf{Beat windows} are 198 samples centred on the annotated R peak, 99 left
and 99 right; windows that would extend past a record boundary are dropped rather
than padded. \textbf{R peaks} come from the WFDB \texttt{atr} annotations, never
from a detector, so detector error is excluded from this study by construction.
\textbf{Labels} follow the AAMI N/S/V/F/Q mapping; non-beat annotation symbols
and symbols outside the mapping are skipped.
\textbf{RR features} are the preceding and following intervals in seconds,
normalized with MIT--BIH DS1 statistics only ($\mu = 0.7785$\,s,
$\sigma = 0.4956$\,s) and clipped to $\pm 2$; no target-database statistic ever
enters normalization.

\subsection{Replay quantization, stated exactly}
\label{sec:quant}

Replay values are stored INT8 with \textbf{symmetric, per-tensor} quantization.
For a stored tensor $x$ the buffer holds
\begin{equation}
q=\operatorname{clip}\!\left(\operatorname{round}\!\left(x/s\right),-127,127\right),
\qquad
\hat{x}=s\,q,
\qquad
s=\frac{\max|x|}{127},
\label{eq:quant}
\end{equation}
where the scale $s$ is calibrated by absolute maximum over the whole buffer at
the moment it is written, and the zero point is fixed at $0$ by symmetry. One
FP32 scale is stored per quantized tensor: one for a post-pool buffer, and two
for a raw buffer, which quantizes the beat and the RR pair on separate scales.
Per-row scales are deliberately \emph{not} used, because a scale per exemplar
would add 4\,B per entry and is already charged as buffer metadata in
Eq.~\eqref{eq:completebudget2}. During adaptation the buffer is dequantized back
to FP32 by Eq.~\eqref{eq:quant} before entering the graph, so training sees
exactly the INT8-rounded values a device would replay --- the quantization error
is inside the experiment, not idealized away.

\subsection{Replay-buffer construction}
\label{sec:replay_policy}

The replay policy is held fixed across every arm so that the comparison isolates
allocation rather than selection quality. The buffer is drawn once from the DS1
source pool by a \textbf{maximum-volume class-balanced selector}: each of the
five AAMI classes receives an equal quota $\lfloor N/5 \rfloor$, and because rare
classes exhaust their pool before the quota is met, the unfilled slots are
redistributed round-robin among classes that still have examples rather than
being wasted. The selector is seeded by the run seed, so the buffer is
\emph{fixed across patients} within a seed and \emph{resampled across seeds};
patient-to-patient variation therefore cannot come from a changing buffer, while
seed-to-seed variation includes buffer variation, which is the conservative
choice.

Values are stored INT8 with a per-tensor scale and zero point. A raw entry stores
the ECG window \emph{and} its RR pair, since the model consumes both and a window
alone cannot be replayed; the honest per-item cost is therefore 200\,B of values
rather than the 198\,B of the window alone, and the budget assertion uses 200.
Post-pool entries are extracted once from the frozen source encoder and are
therefore \emph{stale by construction}: they are never re-extracted as adaptation
proceeds, which is exactly why a post-pool buffer cannot track a representation
that is itself changing --- and is one reason a post-pool buffer pairs naturally
with a frozen encoder. Raw entries carry no such staleness because they are
re-forwarded through the current model on every update, at a greater but
unmeasured compute cost (Appendix~\ref{sec:embedded}). Replay and target beats are
concatenated into one training set rather than mixed at a fixed per-batch ratio,
so the effective replay ratio is $N_{\mathrm{replay}} / (N_{\mathrm{replay}} +
N_{\mathrm{target}})$. There is no eviction policy: the buffer is written once
and never replaced, since the source distribution is static.

\subsection{Metrics and statistical analysis plan}

\paragraph{The primary metric, stated exactly.} For patient $p$ let $n_{p,c}$ be
the number of test beats of AAMI class $c$. The scored class set and the primary
metric are
\begin{equation}
\mathcal{C}_p = \{c \in \{\mathrm{N},\mathrm{S},\mathrm{V}\} : n_{p,c} \ge n_{\min}\},
\qquad
F1_{\mathrm{macro},p} = \frac{1}{|\mathcal{C}_p|}\sum_{c \in \mathcal{C}_p} F1_{p,c},
\label{eq:metric}
\end{equation}
with $n_{\min} = 1$: a class counts as present if the patient's test segment
contains at least one beat of it. The remaining conventions are fixed as follows.
The class set is N/S/V. Fusion (F) is excluded from the primary metric because it
is absent from most DS2 test segments, so including it would make the
denominator $|\mathcal{C}_p|$ vary for reasons unrelated to performance; F is
reported separately as a pooled secondary metric, where the encoder arms show
their largest relative gain. The unclassifiable class Q is excluded because the
source model is not trained to separate paced and unclassifiable beats and the
protocol makes no claim about them. A class with $n_{p,c} = 0$ is omitted from
the average rather than scored as zero, since scoring an absent class as zero
would penalize a patient for the composition of their own record. Precision and
recall use the zero-division convention $0/0 = 0$, which applies only when a
present class receives no predictions.

Secondary metrics are pooled accuracy, per-class S/V/F sensitivity, precision,
and F1, the fraction of patients improved versus A0, and source-domain \macrof{}
change $\Delta^{\mathrm{src}}$ (Eq.~\ref{eq:gain_and_srcchange}), which is negative when
adaptation costs source performance. Pooled and per-patient summaries are both reported because a
single record (232) contributes a large fraction of the pooled S beats, so pooled
scores alone would be misleading. Because the improved-patient count is sensitive
to how small changes are treated, we report it both threshold-free and at a
meaningful-change threshold of $0.02$ \macrof{} (Table~\ref{tab:harm}).

The primary metric, seed count, paired tests, effect-size definition, and Holm
comparison family were fixed before the final runs. Record 202 was excluded after
a subject-identity audit showed that it belongs to the same patient as DS1
record 201 (Section~\ref{sec:cohort}); all affected analyses were rerun without
changing the methods or the comparison family.

Seeds are averaged per patient before
patient-wise paired tests. For each planned comparison we report the mean and
median paired difference, a \textbf{patient bootstrap} 95\% confidence
interval ($10{,}000$ resamples), a paired Wilcoxon signed-rank $p$-value,
rank-biserial effect size, and a Holm correction across the six-comparison
family $\{$A4--A1, A5--A1, A4--A2, A4--A3, A4--A0, A5--A0$\}$. A leakage and
reproducibility audit (no DS1/DS2 record overlap, no adaptation/test window
crossing, no \emph{primary-cohort} test label used for patient inclusion,
replay selection, early stopping, hyperparameter selection, or adaptation ---
external inclusion and the class-aware oracle deliberately do use test-window
labels and are excluded from this check (Sections~\ref{sec:rq6}
and~\ref{sec:rq7}) --- source-only
normalization statistics, stored seeds and library versions, recorded source
checkpoint hash) accompanies the run and returns PASS\,11 / FIXED\,1 / ASSESSED\,1 / WARN\,1 /
FAIL\,0.

\paragraph{The bootstrap, stated exactly.} Every interval in this paper is a
\emph{patient bootstrap}, and we name it that way consistently because seeds are
collapsed before resampling: (i) average the five seeds within each patient to
give one score per patient; (ii) resample the 21 patients with replacement;
(iii) recompute the arm mean or the paired difference on the resampled patients;
(iv) repeat $10{,}000$ times and take the $2.5$th and $97.5$th percentiles.
Seeds are therefore not resampled, and we \textbf{do not} call this a
hierarchical bootstrap --- that term is reserved for a procedure that resamples
seeds inside every resampled patient, which would be the right choice if
seed variance were the quantity of interest. External-database intervals use the
same procedure with subjects in place of patients.
